\section{Continuous coefficients and full-group duality}
\label{sec:continuous}

The coefficient modules below include compact formal lattices, their
localizations, and spaces of convergent Laurent series.  We first specify
one cochain convention that applies to all three. Finite free resolution
terms make the cohomology of compact coefficients compact Hausdorff;
this is the input used when countability is upgraded to finiteness.
The marked Laurent support module needs a second conclusion: duality
between its group homology and cohomology for the full ordinary
stabilizer. For that module, trace-one elements for finite subgroups
replace averaging by their orders and permit the full-group duality
argument below. The two conclusions have different hypotheses; only
the second requires the trace condition. In this section $\kk$ is a
finite field of characteristic $p$, and $G$ is a compact $p$-adic
analytic group.

\begin{definition}[Complete steps]
\label[definition]{def:cts-steps}
A coefficient module with complete steps is a vector space
$M=\bigcup_{a\geq 0}M_a$, where the $M_a$ are complete Hausdorff linearly
topologized $\kk$-spaces and the inclusions are continuous injections.
The inclusions need not induce the topology on the smaller space.  An
action of $G$ is required to have the following property: for every $a$
there is a $b$ for which $G M_a\subset M_b$ and
$G\times M_a\longrightarrow M_b$ is continuous.  Morphisms satisfy the
analogous condition with no group variable.  We use the bounded-step
continuous cochains
\[
 C^j_{\cts}(G,M)=\colim_a C_{\cts}(G^j,M_a)
\]
with the usual inhomogeneous differential.  Their cohomology is denoted
$H^j(G,M)$.  The same convention applies to continuous functions on any
compact profinite parameter space $\mathcal D$: its complete steps are
$C_{\cts}(\mathcal D,M_a)$ with the topology of uniform convergence.
\end{definition}

All subsequent assertions about these coefficient modules refer to this
specified complex; they do not replace an algebraic image by its closure.
Finite-dimensional discrete modules and compact modules are special
cases.  A localization of a finite module over
$\kk[[x,y_1,\ldots,y_s]]$ has the steps $x^{-a}M_0$, with their adic
topologies.  Laurent modules have the complete steps specified in
\cref{sec:laurent}.

\begin{lemma}[Exact cochains]
\label[lemma]{lem:cts-exact}
Consider an exact sequence of underlying vector spaces
\[
 0\longrightarrow M'\xrightarrow{\iota}M\xrightarrow{q}M''
 \longrightarrow0
\]
whose maps and actions respect complete steps.  Suppose that:
\begin{enumerate}[label=\textup{(\roman*)}]
\item a bounded-step continuous map from a compact profinite space to
      $M$ with image in $\iota(M')$ is bounded-step continuous as a map to
      $M'$;
\item for every compact subset $T$ of a step of $M''$, there is a
      continuous map $s:T\to M_b$ into one step of $M$, with $qs=\id_T$.
\end{enumerate}
Then the corresponding sequence of continuous cochain complexes is
short exact.  In particular, it has the usual long exact cohomology
sequence.  These hypotheses hold for strict short exact sequences of
compact $\kk$-spaces, for short exact sequences of finite modules over a
complete Noetherian local $\kk$-algebra with finite residue field, and
for their localizations with the lattice steps just described.
\end{lemma}

\begin{proof}
The kernel assertion is (i).  A cochain in the quotient has compact
image $T$ in a single step, and composing it with the section in (ii)
lifts that cochain.  The section is not required to commute with $G$.
Thus every degree of the cochain sequence is exact.

For compact vector spaces, continuous duality with discrete vector
spaces turns a strict surjection into an injection.  Splitting that
injection as a map of vector spaces and dualizing gives a continuous
$\kk$-linear section.  In the Noetherian local case the modules are
compact.  A submodule is closed, and its induced topology is its adic
topology by Artin--Rees.  This proves (i), while the compact splitting
proves (ii).  For a localization, a fixed quotient lattice is the image
of a compact source lattice: choose lifts of finitely many module
generators and a common denominator.  Its map from that source lattice
is a strict compact surjection and has a continuous linear section.
The kernel in a fixed lattice is again a finite module, so Artin--Rees
gives (i) after a common enlargement of the denominator.  This also
proves the assertions for finite twists and finite direct sums.
\end{proof}

\begin{lemma}[Compact envelopes and comparison]
\label[lemma]{lem:cts-comparison}
Let $M$ satisfy \cref{def:cts-steps}.
\begin{enumerate}[label=\textup{(\roman*)}]
\item Every compact subset of a step of $M$ is contained in a compact
      $G$-stable $\kk$-subspace $N$ lying in one complete step.  The
      action on $N$ is continuous.  These $N$ form a filtered system
      whose union is $M$, and
      \[
       C^\bullet_{\cts}(G,M)
       =\colim_N C^\bullet_{\cts}(G,N).
      \]
\item The ring $\kk[[G]]$ is Noetherian, and the trivial compact module
      $\kk$ has a resolution $P_\bullet$ by finitely generated free
      compact $\kk[[G]]$-modules.  It may be unbounded.  The comparison
      map
      \[
       \Hom_{\kk[[G]],\cts}(P_\bullet,M)
          \simeq C^\bullet_{\cts}(G,M)
      \]
      is a natural quasi-isomorphism.  Every term on the left is a
      finite power of $M$.
\item A right resolution by finitely generated free compact modules
      similarly defines
      $H_i(G,M)$.  For compact $M$, cohomology and homology computed
      this way are compact Hausdorff, and continuous duality
      interchanges homology and cohomology, with the contragredient
      action on $M^\vee=\Hom_{\cts}(M,\kk)$.
\end{enumerate}
For $G=\Zp$ with generator $\gamma$, the cochain complex is
$[M\xrightarrow{\gamma-1}M]$ in degrees $0,1$.
\end{lemma}

\begin{proof}
Let $T\subset M_a$ be compact.  The set $GT$ is compact in some $M_b$.
Let $N$ be its closed linear span there.  Modulo any open linear
subspace of $M_b$, the image of $GT$ is finite.  Since $\kk$ is finite,
its linear span in that discrete quotient is finite as well.  Thus
$N$ is totally bounded and complete, hence compact.  For any $g\in G$,
the action map $M_b\to M_c$ is continuous after enlarging $c$ uniformly
in $g$.  The image of $N$ in the Hausdorff space $M_c$ is compact and
closed.  It contains the image of the linear span of $GT$, and
continuity therefore shows $gN\subset N$.  The injection $N\to M_c$
is a homeomorphism onto its image, so the restricted action is
continuous.  A finite sum of such spaces is again compact in a common
step.  Every element, and every compact cochain image, is contained in
one of them.  A map from a compact domain whose image lies in $N$ is
continuous into $N$, by its subspace topology in $M_b$.  This proves
(i), including the assertion about cochains.

Choose a normal uniform open subgroup $K\subset G$.  For the
augmentation filtration of $\kk[[K]]$, the associated graded ring is
a polynomial ring on $\dim G$ variables.  Thus $\kk[[K]]$ is
Noetherian, and the finite crossed-product extension $\kk[[G]]$ is
Noetherian as well.  Successively resolving the finitely generated
kernels of a finite free presentation gives $P_\bullet$.  For a
compact module $N$, compare it with the compact bar resolution.
Its terms are induced from compact $\kk$-spaces and are projective:
every compact $\kk$-space is projective by continuous linear duality,
and completed induction preserves that lifting property.
Both are projective resolutions, and their augmentations and
acyclic kernels split as compact $\kk$-spaces by continuous duality.
The lifting construction for resolutions gives chain maps and chain
homotopies in the compact category.  Applying continuous Hom gives
the usual continuous bar cochains.  For comparison with the
ind-profinite formulation under its countability hypotheses, see
\citet[Proposition~8.2 and Corollary~8.4]{boggiCorobCook}.
The compact splitting argument just given does not impose
second countability on $N$ and proves that all its comparison
homotopies are continuous.

There is a well-defined $\kk[[G]]$-action on $M$: integrate on a compact
$G$-stable envelope of a given element.  It is independent of that
envelope.  Since every $P_j$ is finite free, its Hom into $M$ is the
filtered union of its Hom into the $N$.  Filtered colimits of vector
spaces are exact.  Passing the compact comparison and its homotopies
through this colimit proves (ii); in particular no new completion is
introduced.  The same argument applies to the finite free right
resolution and tensor products.  For compact $N$, every differential
has compact, hence closed, image in a finite power of $N$.
Continuous duality is exact on compact vector spaces and takes these
finite-power complexes to the corresponding dual complexes.  This
proves (iii).  Finally,
$\kk[[\Zp]]=\kk[[\gamma-1]]$, and multiplication by $\gamma-1$ gives
the stated two-term free resolution.
\end{proof}

We next make the finite-isotropy hypothesis explicit.  Its verification
for the geometric coefficient algebras will be given in
\cref{prop:mark-finite-isotropy}.

\newpage
\begin{definition}[Trace elements]
\label[definition]{def:cts-trace}
Let $A$ be a commutative $\kk$-algebra with a $G$-action.  A semilinear
$A$-module $M$ with complete steps is trace-admissible if, for every
finite subgroup $F\subset G$, there is $t_F\in A$ satisfying
\[
       \sum_{f\in F}f(t_F)=1,
\]
and multiplication by $t_F$ carries each complete step continuously
into some complete step.  The action is semilinear in the sense that
$g(am)=g(a)g(m)$.
\end{definition}

We will need an elementary local fact about the finite stabilizers of
profinite spaces.  The characteristic-zero Lie group in its proof is
the group $K$, even though the coefficient field has characteristic
$p$.

\begin{lemma}[Equivariant local sections]
\label[lemma]{lem:cts-local-section}
Let $K$ be a normal uniform open subgroup of $G$, and let $X$ be a profinite
$G$-space on which $K$ acts freely.  Suppose that $X\to Y=K\backslash X$
has a continuous $K$-equivariant trivialization as a principal
$K$-space.  Put $H=G/K$.  If $x\in X$ and $y$ is its image, the finite
group $G_x$ maps isomorphically to $H_y$.  There is an $H_y$-invariant
clopen neighborhood $V$ of $y$ and a section $s:V\to X$, with $s(y)=x$,
which is equivariant for $G_x\simeq H_y$.
\end{lemma}

\begin{proof}
An element of $H_y$ has a lift taking $x$ into its $K$-orbit.  Correcting
that lift by the unique element of $K$ gives an element fixing $x$.
Since $K$ acts freely, the resulting homomorphism $G_x\to H_y$ is
bijective.  Write $F=G_x$.  A continuous section with value $x$ at $y$
exists by the assumed trivialization.  Its failure to be $F$-equivariant
is a continuous nonabelian cocycle
$c_f\in C(V,K)$ for the action combining conjugation by $F$ and its
action on $V$.  All $c_f(y)$ equal $1$.  Shrinking to an $F$-invariant
clopen neighborhood makes all these cocycles arbitrarily close to $1$.

Here is the required small-cocycle argument, including the case
$p\mid |F|$.  Choose an $F$-stable sufficiently small Lie lattice in
$\Lie(K)$ on which exponential, logarithm, and the
Campbell--Hausdorff formula converge.  In the Banach space of continuous
functions on $V$ with values in $\Lie(K)$, use the supremum norm of an
$F$-invariant norm.  Put $\lambda_f=\log c_f$ and
$\epsilon=\max_f\|\lambda_f\|$.  Convergence of the
Campbell--Hausdorff series gives, for a fixed constant $c$, the bound
\[
 \|\lambda_{fg}-\lambda_f-f\lambda_g\|\leq c\epsilon^2.
\]
With $b=|F|^{-1}\sum_f\lambda_f$, summing this identity over $g$
gives
\[
 \|\lambda_f+fb-b\|\leq c|\,|F|\,|_p^{-1}\epsilon^2,
 \qquad \|b\|\leq |\,|F|\,|_p^{-1}\epsilon.
\]
Changing the section by $\exp(-b)$ changes its cocycle to
$\exp(-b)c_f f(\exp b)$.  Another application of the same convergence
estimate bounds the new logarithms by $c'\epsilon^2$, for a constant
$c'$ depending only on $F$ and the chosen Lie lattice.  Choose the
initial neighborhood small enough that these exponentials stay in
the convergence lattice and $c'\epsilon<1$.  Iteration then gives
uniformly convergent continuous changes of section, since the
successive logarithms tend to zero at least quadratically.  The limit
has trivial cocycle.  Every change is the identity at $y$, so the
limiting section still takes $y$ to $x$.  Division by $|F|$ was only
performed in $\Qp$; its fixed norm loss was absorbed in the initial
choice of neighborhood.
\end{proof}

\begin{proposition}[A uniform finite-isotropy bound]
\label[proposition]{prop:cts-isotropy-bound}
Fix a normal uniform open subgroup $K\subset G$, and put
$d_G=\dim G$ and $a_G=[G:K]$.  If $M$ is trace-admissible, then
\[
                 H^j(G,M)=0\qquad(j>a_Gd_G).
\]
The same bound holds for $M\otimes_\kk V$ for every finite-dimensional
continuous $G$-representation $V$.  Neither the bound nor the choice
of $K$ depends on $M$ or $V$.
\end{proposition}

\begin{proof}
First consider a finite subgroup $F$ and a semilinear module $L$ with
a trace element $t=t_F$.  The map
\[
 L\longrightarrow \kk[F]\otimes_\kk L,
 \qquad v\longmapsto\sum_{f\in F}f\otimes t f^{-1}v
\]
is $F$-linear when $F$ acts on the first factor of the target.  Its
composite with $f\otimes v\mapsto fv$ is multiplication by
$\sum_f f(t)=1$.  Thus $L$ is a direct summand of an induced module.
The same formulas on continuous cochains, which are finite sums of
bounded continuous operators, show $H^{>0}(F,L)=0$ in our convention.
They apply also to continuous function modules with pointwise
multiplication by $t$, and to any finite-dimensional twist.

The trivial $\kk[[K]]$-module has a finite free resolution of length
$d_G$.  For finite cohomological dimension and uniform-group Poincare
duality, see \citet[Chapter~V, Proposition~2.5.7.1 and
Theorem~2.5.8]{lazard1965} and \citet[Lemma~4.15 and
Proposition~4.16]{bghs}.  Thus $\kk$ has projective dimension $d_G$,
and the Noetherian local ring $\kk[[K]]$ has finitely generated
free terms in a minimal resolution.  A sufficiently deep uniform
subgroup is used also when $p=2$.
Tensor-induce this resolution from $K$ to $G$:
take its completed tensor product over the $a_G$ cosets and let $G$
permute the factors, using the usual signs of complexes and the
conjugate $K$-actions.  Denote the resulting augmented complex by
$T_\bullet\to\kk$.  It is concentrated in degrees $0$ through
$a_Gd_G$.  The original augmented resolution is split exact as a
complex of compact $\kk$-spaces.  Its tensor product is therefore
again a split resolution of $\kk$ in that category.  This is the
tensor-induction construction of
\citet[Theorem~5.2]{symonds2007}.

Each term of $T_\bullet$ is a finite sum of signed permutation modules
$\kk[[X]]$.  Here $X$ is a finite disjoint union of products of copies
of $K$, one for each tensor factor.  The diagonal $K$-action is free,
and setting the first coordinate equal to $1$ gives a continuous
trivialization of $X\to K\backslash X$.  A possible sign is a
locally constant rank-one coefficient and will be carried along.
We claim that
\begin{equation}
 H^j\bigl(G,\Hom_{\kk,\cts}(\kk[[X]],M)\bigr)=0\quad(j>0).
 \label{eq:cts-permutation-acyclic}
\end{equation}
The Hom module is $C_{\cts}(X,M)$, with complete uniform-convergence
steps.  A $K$-trivialization identifies it, after untwisting the
diagonal action on $M$, with a continuously coinduced $K$-module.
Evaluation in the first coordinate contracts the coinduced bar
complex, so its positive $K$-cohomology vanishes.  All operators in
that contraction are bounded continuous operators on function
steps.  The Hochschild--Serre double complex therefore reduces its
$G$-cohomology to the cohomology of $H=G/K$ on the $K$-invariants.

Apply \cref{lem:cts-local-section} at a point of $Y=K\backslash X$.
After shrinking its neighborhood $V$ further, its translates by
elements outside $H_y$ are disjoint: first separate $y$ from the
finitely many distinct points $hy$, and then intersect the resulting
clopen neighborhoods.  On the induced neighborhood $HV$, the
$K$-invariant functions are induced from $H_y\simeq G_x$ acting on
$C_{\cts}(V,M)$, with its possible rank-one sign.  The equivariant
section makes this the ordinary semilinear action on values together
with the action on $V$.  Its higher $H_y$-cohomology vanishes by the
trace-element splitting.  These induced neighborhoods cover the
compact space $Y$.  Choose a finite such cover and take successive
differences of its $H$-invariant clopen members to obtain a disjoint
finite invariant clopen partition.  Each piece retains its local
description and splitting.  Finite additivity and finite-group
Shapiro now prove \eqref{eq:cts-permutation-acyclic}.

Finally form the first-quadrant double complex
\[
 C^u_{\cts}\bigl(G,\Hom_{\kk,\cts}(T_v,M)\bigr).
\]
The compact $\kk$-linear contracting homotopy of the augmented
$T_\bullet$ acts on these Hom modules: a continuous map from a
compact domain into a step has image in a compact envelope, so
integration and the contracting homotopy take place in complete
steps by \cref{lem:cts-comparison}.  Computing first in the $v$
direction identifies the total cohomology with $H^*(G,M)$.
Computing first in the $u$ direction and using
\eqref{eq:cts-permutation-acyclic} leaves only $u=0$ and
$0\leq v\leq a_Gd_G$.  This proves the bound.  Tensoring $M$ with
$V$ preserves every step condition and the trace elements, so the
same construction proves the last assertion with the identical
bound.
\end{proof}

The preceding bound is deliberately larger than the dimension.
It is its uniformity under finite twists that permits passage from
torsion-free open subgroups to the full group, including at primes
where $G$ has torsion.

\begin{proposition}[Full-group duality]
\label[proposition]{prop:cts-full-duality}
Assume that the orientation character of $G$ is trivial over $\kk$.
For every trace-admissible $M$ there are natural isomorphisms
\[
                  H_i(G,M)\simeq H^{d_G-i}(G,M)
                  \qquad(i\in\Z).
\]
Here homology is zero in negative degrees, and cohomology is zero in
negative degrees.  In particular $H^j(G,M)=0$ for $j>d_G$.
The ordinary stabilizer $P_n$ has trivial orientation character and
$d_G=n^2$, so this assertion applies to it without a twist.
\end{proposition}

\begin{proof}
Write $R=\kk[[G]]$, $R_K=\kk[[K]]$, $H=G/K$, and
$\Sigma=\kk[H]$.  Regard $\Sigma$ as the quotient bimodule of $R$.
It is perfect as a left or right $R$-module, since it is induced
from the length-$d_G$ free resolution over $R_K$.  Uniform-group
duality, with its conjugation action retained, gives
\begin{equation}
 \RHom_R(\Sigma,R)\simeq\Sigma[-d_G]
 \label{eq:cts-bimodule-duality}
\end{equation}
as a $(\Sigma,R)$-bimodule.  To recall the equivariance in this
formula, restrict $R$ to $R_K$ and decompose it into its finitely
many cosets.  The dual of the uniform resolution has cohomology
only in degree $d_G$, with one orientation line on each coset.
Right translation permutes those lines by the right regular action;
precomposition by a quotient element gives the left regular action
and its action on the orientation line.  The latter is trivial by
hypothesis.  This proves exactly the two commuting actions asserted
in \eqref{eq:cts-bimodule-duality}.  The uniform duality, orientation
identification, and compatibility with conjugation are also recorded in
\citet[Propositions~4.16 and~4.36, Theorem~4.19 and
Remark~4.23]{bghs}.

Put
\[
 C=\RHom_R(\Sigma,M)=\RGamma(K,M),\qquad
 B=\Sigma\otimes_R^{\mathbf L}M.
\]
Perfectness permits evaluation of the dual resolution against $M$,
so \eqref{eq:cts-bimodule-duality} gives a natural equivalence
$C\simeq B[-d_G]$ in the derived category of $\Sigma$-modules.
The cohomology of $C$ lies in degrees $0,\ldots,d_G$.  The finite
free comparisons in \cref{lem:cts-comparison} show that these
derived expressions compute precisely the specified continuous
cochains and homology, even when $M$ is a union of steps.

We prove that the finite-group norm
\begin{equation}
       \kk\otimes_\Sigma^{\mathbf L}C
           \longrightarrow\RHom_\Sigma(\kk,C)
       \label{eq:cts-finite-norm}
\end{equation}
is an equivalence.  This is the point where finite subgroup
cohomology cannot be discarded.  For any finite-dimensional left
$\Sigma$-module $V$, derived adjunction gives
\[
 \Ext^j_\Sigma(V,C)
       =H^j(G,V^*\otimes_\kk M).
\]
By \cref{prop:cts-isotropy-bound} these groups vanish for
$j>a_Gd_G$, uniformly in $V$.  Choose a bounded-below injective
resolution $J^\bullet$ of $C$.  If $N$ is greater than both $d_G$
and $a_Gd_G$, dimension shifting along its exact tail gives
\[
 \Ext^1_\Sigma(V,\ker(J^N\to J^{N+1}))
       =\Ext^{N+1}_\Sigma(V,C)=0.
\]
Every quotient $\Sigma/I$ by a left ideal is finite-dimensional.
Baer's criterion therefore says that this last kernel is injective.
Truncating there gives a bounded complex of injective modules
representing $C$.

For a finite group algebra, every injective module, including an
infinite-dimensional one, is projective.  Indeed, the canonical
injection of any module $L$ into
$\Hom_\kk(\Sigma,L)$ splits if $L$ is injective; the target is free
because the symmetric pairing on $\kk[H]$ identifies coinduction
with induction.  On a free module of arbitrary rank the ordinary
norm is an isomorphism from coinvariants to invariants: on each
copy of $\kk[H]$ it sends the class of $1$ to $\sum_{h\in H}h$.
It is consequently an isomorphism on projective summands.  Applying
this degreewise to the bounded injective-projective model of $C$
proves \eqref{eq:cts-finite-norm}, with no infinite totalization or
completion.

The left-hand side of \eqref{eq:cts-finite-norm}, shifted by $d_G$,
is
\[
 (\kk\otimes_\Sigma^{\mathbf L}C)[d_G]
       \simeq\kk\otimes_R^{\mathbf L}M;
\]
the right-hand side is $\RGamma(G,M)$.  Taking cohomology in
degree $-i$ of the shifted equivalence gives
$H_i(G,M)=H^{d_G-i}(G,M)$, with the stated grading.

For $G=P_n$, its Lie algebra is the central division algebra
$D_n$ of degree $n$ over $\Qp$.  Conjugation by $a\in D_n^\times$
has determinant
\[
 \det(L_a)\det(R_a)^{-1}
   =\operatorname{Nrd}(a)^n\operatorname{Nrd}(a)^{-n}=1.
\]
The orientation character is the reduction of this determinant
character, hence is trivial.  The dimension of $D_n$ is $n^2$.
\end{proof}

\begin{remark}
For a nontrivial orientation character, the same proof retains its
one-dimensional line in \eqref{eq:cts-bimodule-duality} and gives
the corresponding orientation-twisted duality.  We only use the
untwisted case stated above.  The proposition makes no assertion
that arbitrary $\kk[[G]]$-modules have finite cohomological
dimension: its trace-element hypothesis is essential when $G$
contains $p$-torsion.
\end{remark}
